\worksheettitle
  {Домашняя практика}
  {\modulemark{M07-H} Равные промежутки}

\studentnameblock

\begin{practicebox}[1. Число промежутков]
  \begin{enumerate}[label=\alph*)]
    \item между 6-й и 15-й точками: \answerblank[18mm];
    \item между 18-й и 42-й: \answerblank[18mm];
    \item $12$ подряд стоящих столбов образуют \answerblank[18mm]
      промежутков.
  \end{enumerate}
\end{practicebox}

\begin{practicebox}[2. Найдите расстояние]
  Точки отмечены через $6$ мм.
  \begin{enumerate}[label=\alph*)]
    \item между 4-й и 13-й точками: \answerblank[55mm];
    \item между 25-й и 37-й: \answerblank[55mm].
  \end{enumerate}
\end{practicebox}

\begin{practicebox}[3. Обратные задачи]
  \begin{enumerate}[label=\alph*)]
    \item Между 7-й и 12-й точками $35$ см. Найдите шаг.
    \item От 20-й точки вправо прошли $48$ см при шаге $6$ см.
      Найдите номер конечной точки.
  \end{enumerate}
  \writinglines{3}
\end{practicebox}

\begin{practicebox}[4. Ряд объектов]
  От первого до последнего фонаря $45$ м. Фонари стоят через $5$ м и есть
  в обоих концах. Сколько всего фонарей? \writinglines{2}
\end{practicebox}

\begin{warningbox}[5. Исправьте ошибку]
  «Между 8-й и 20-й точками $20-8+1=13$ промежутков».
  Назовите причину и запишите правильное решение. \writinglines{2}
\end{warningbox}
